\section{The original raw density in a complete local norm}
\label{sec:v88-raw}
The surface theorem applies to every first-exit component, including
the zero-residence tail. Combined with the retained fractional source,
it gives a norm estimate for the original bad density itself. The
roof density is not averaged in the conclusion.

\begin{lemma}[A truncated potential in a local norm]
\label{lem:v88-potential-norm}
Let $0<\kappa<1$, $q>1$, and $1/q>\kappa$. For a nonnegative
locally integrable $v$, put $P_rv(t)=\int_0^r s^{-\kappa}v(t+s)\,ds$.
Then
\begin{equation}\label{eq:v88-potential-norm}
 \mathcal N_q(P_rv;J)\le\|P_rv\|_{L^q(J)}
 \le (1-\kappa q)^{-1/q}r^{1/q-\kappa}
                                    \int_{J+(0,r)}v(u)\,du.
\end{equation}
Also, if $1\le p<q$,
\begin{equation}\label{eq:v88-embedding}
 \|f\|_{L^p(J)}\le
   \left(\frac q{q-p}\right)^{1/p}
         |J|^{1/p-1/q}\mathcal N_q(f;J).
\end{equation}
\end{lemma}
\begin{proof}
For fixed $u$, the kernel $\one_{\{0<u-t<r\}}(u-t)^{-\kappa}$
has $L^q(dt)$ norm at most
$(1-\kappa q)^{-1/q}r^{1/q-\kappa}$. Minkowski's integral
inequality over $u\in J+(0,r)$ gives the second inequality in
\eqref{eq:v88-potential-norm}; H\"older on each measurable $E\subset J$
gives the first. Apply this to truncated kernels and pass monotonically
if needed. To prove \eqref{eq:v88-embedding}, the definition of
$\mathcal N_q$ implies
$|\{|f|>\lambda\}\cap J|\le\min\{|J|,(\mathcal N_q/\lambda)^q\}$.
Integrate the distribution function against $p\lambda^{p-1}$ and
split at $\lambda=\mathcal N_q|J|^{-1/q}$.
\end{proof}

\begin{theorem}[Complete original bad-source control]
\label{thm:v88-bad-source}
For every $1<q<32/31$ and fixed $L>0$, under the explicitly
retained inputs,
\begin{equation}\label{eq:v88-bad-source}
 \limsup_m\sup_{R,n,k}\sup_{|J|=L}
         \mathcal N_q(m^2b^\varepsilon;J)
                           \le C_q L^{1/q}\varepsilon^{1/16}
\end{equation}
for every fixed sufficiently small $\varepsilon>0$. All initial
momenta, positive total residences, collision incidences, terminal
tests and exact labels are retained. No original source restriction
is removed in forming $b^\varepsilon$.
\end{theorem}
\begin{proof}
The source-shift inequality \eqref{eq:v83-shift-input} and the
complete two-sided optical payment \eqref{eq:v88-optical-input}
give, for each fixed $r>0$,
\[
 b^\varepsilon\le 2A_rb^\varepsilon
       +C\overline{\mathcal F}_r+C_\eta\mathsf P_{\eta,r}v^\eta
       +C_\Theta(1+\varepsilon^{-1})\sqrt r\,A_rp.
\]
Every term is a complete source term, not a sum with an unknown
physical-word multiplicity. Apply \eqref{eq:v88-weak-norm}, then
take the collision limsup with $\varepsilon,r$ and all auxiliary
bands fixed. The fixed-width source bounds in
Input~\ref{input:v83-local} give
\[
 \limsup_m\sup_{R,n,k}m^2\|A_rb^\varepsilon\|_\infty
 \le C\varepsilon^{1/16},\qquad
 \limsup_m\sup_{R,n,k}m^2\|A_rp\|_\infty\le C.
\]
Theorem~\ref{thm:v88-full-surface} controls the surface by
$C_{q,L}r^{\nu_q}$. Choose $\eta=1/4$, so $\kappa=3/8<1/q$;
Lemma~\ref{lem:v88-potential-norm} and the actual fixed-width
fractional input control the potential by
$C_{q,\eta}(L+r)r^{1/q-\kappa}$. Thus the ordered bad-source
norm is at most
\[
 C_qL^{1/q}\varepsilon^{1/16}
 +C_{q,L}r^{\nu_q}
 +C_{q,\eta}(L+r)r^{1/q-\kappa}
 +C_\Theta L^{1/q}(1+\varepsilon^{-1})\sqrt r.
\]
Enlarge auxiliary bands only after the collision limit. Now let
$r\downarrow0$ with $\varepsilon$ fixed. All three powers are
positive. This proves \eqref{eq:v88-bad-source}. Averaging is an
intermediate proof operation; the left side is the raw density.
\end{proof}

\begin{theorem}[Exact-label raw local limit in finite local norms]
\label{thm:v88-raw-limit}
Let
\[
 f_{m,R,n,k}(t)=m^2p_{n,R}(k,m,t)
                    -\mathcal L_{m,R}(k_1,k_2,t,n).
\]
For every fixed $M,L>0$ and every $1<q<32/31$,
\begin{equation}\label{eq:v88-raw-weak}
 \sup_{R,n,k}\sup_{|J|=L}
 \mathcal N_q(f_{m,R,n,k}\one_{\mathcal C_{m,R,n,k}(M)};J)
                                      \longrightarrow0.
\end{equation}
For every $1\le p<32/31$ the corresponding $L^p(J\cap\mathcal C(M))$
norm also tends to zero, uniformly in the same parameters and interval
locations. The reference includes every retained arithmetic peak and
all zero classes; no reference density is divided by.
\end{theorem}
\begin{proof}
For each fixed large $B$ use the original positive remainder
\eqref{eq:v76-remainder}. Theorem~\ref{thm:v88-bad-source} and
monotonicity under restriction to $\mathcal C(M)$ imply
\[
 \limsup_m\sup_{R,n,k}\sup_{|J|=L}
 \mathcal N_q(f_{m,R,n,k}\one_{\mathcal C(M)};J)
                                      \le C_{q,M,L}B^{-1/192}.
\]
Indeed $\varepsilon(B)^{1/16}=A_0^{1/16}B^{-1/192}$, and the
uniform essential bound on $e_B$ contributes at most
$L^{1/q}C_MB^{-1/192}$. The left side is independent of $B$.
Let $B\to\infty$ last. Nonnegativity proves the asserted convergence.
For any stated $p$ choose $q$ strictly between $\max(1,p)$ and
$32/31$ and use \eqref{eq:v88-embedding}.
\end{proof}

The governing objective remains \eqref{eq:v76-target}. The theorem
above is a proved intermediate conclusion for its unchanged raw
error, not a replacement target. Its uniformity is on every fixed
length $L$, not automatically on an interval growing with $m$.
Neither a finite weak norm nor these finite $L^p$ norms control the
normalized essential supremum. The exact original-source criterion
of Supplement Theorem~\ref{supp-v87archive:thm:v87-endpoint} is
therefore retained: vanishing of its critical/zero-residence charge
still has to be proved for the pointwise endpoint. Same-roof ratios
and unconditional likelihood statements are not promoted on the
strength of the local-norm theorem alone.
